\section{Ameaças à validade}
\label{sec:ameacas}

\subsection{Validade de construto}
\label{sec:ameacas-construto}

As métricas observadas são proxies, não medidas diretas de entrega em
produção. Uma release publicada não demonstra que a mudança foi implantada ou
adotada por usuários; uma falha de CI pode bloquear uma mudança antes da
produção e não equivale a uma falha em produção. Assim, os CFRs (a) e (b)
representam eventos operacionais distintos, e nenhum deles mede incidentes
reais.

Essa limitação é especialmente importante para o CFR (b). Na validação de 283
releases em 58 repositórios, com um único avaliador, a heurística obteve
precisão de 0,238 (IC 95\% de Wilson [0,180; 0,308]), recall de 0,870 e
$F_1 = 0,374$, abaixo do limiar de 0,70
(Seção~\ref{sec:validacao-cfrb};
\texttt{analysis/validacao\_cfr\_b/resultados.md}). Das 168 releases marcadas
como corretivas, apenas 40 foram confirmadas; entre as 46 corretivas segundo o
avaliador, seis não foram identificadas. O viés observado é, portanto, de
superestimação do CFR (b), sem que a avaliação de um único julgador permita
separar o erro da heurística de critérios individuais.

A sensibilidade da categoria geral reforça que a escolha operacional afeta a
interpretação: mudaram de categoria 44,4\% dos 259 repositórios comparáveis
entre C1 e C2 ($\kappa = 0,311$), 28,5\% dos 260 entre C1 e C3
($\kappa = 0,625$) e 41,0\% dos 261 entre C2 e C3 ($\kappa = 0,428$)
(\texttt{analysis/rq07/rq07\_resultados.md}). Como C2 usa o CFR (b), essas
mudanças não isolam o efeito da definição de entrega do possível erro de sua
heurística. Não se deve interpretar essas proporções como mudanças de
desempenho em produção.

\subsection{Validade interna}
\label{sec:ameacas-interna}

As associações e diferenças entre grupos são observacionais e não controlam
todos os fatores que podem influenciar as métricas. Idade e atividade dos
repositórios, por exemplo, podem afetar simultaneamente volume de releases,
commits e execuções de CI. A RQ06 compara linguagem, quartis de estrelas e
quartis de contribuidores, mas não ajusta por idade ou atividade; seus dois
resultados significativos entre 12 testes dizem respeito ao lead time e não
estabelecem efeitos causais (\texttt{analysis/rq06/rq06\_resultados.md}).

O CFR (b) também depende de uma heurística: uma alteração apenas de \textit{patch}
ou mensagens contendo termos como \texttt{fix}, \texttt{hotfix} ou
\texttt{revert} podem classificar uma release como corretiva sem evidência de
que a entrega anterior falhou. A validação de avaliador único e seu
$F_1 = 0,374$ limitam a força das conclusões que usam essa variável. Uma
mitigação é rotular releases de forma independente por múltiplos avaliadores,
resolver divergências por consenso e reavaliar a heurística em uma amostra
separada antes de interpretar o CFR (b).

O lead time usa \texttt{commit.author.date}, que pode não representar o momento
em que o commit entrou no branch padrão: rebase e \textit{squash merge} podem
preservar ou alterar datas de autoria e modificar a relação temporal entre
commits e releases. Essa fonte de incerteza deve ser considerada junto à
variabilidade das datas de publicação. Como mitigação, futuras análises podem
comparar a definição baseada na data de autoria com a data de integração no
branch padrão, documentando a regra e reportando a sensibilidade dos
resultados; a análise atual não estima esse efeito separadamente.

Há ainda censura à direita na janela de observação. No CFR (b), 251 das 8.518
releases (2,9\%) não dispõem dos sete dias seguintes completos; na recuperação,
75 dos 4.119 episódios (1,8\%) permanecem sem encerramento
(\texttt{analysis/rq03\_rq04/rq03\_rq04\_resultados.md}, seção
\textit{Censura}). Esses eventos foram mantidos como censurados, sem imputar
tempos ou resultados. Ainda assim, a janela pode subcontar correções e
recuperações tardias. Uma mitigação é estender a observação por ao menos sete
dias após a janela principal para o CFR (b), reportar separadamente os
episódios censurados e repetir a análise com uma janela de acompanhamento
adicional.
